% Pursuing Stacks (cotes 134-2 to 134-8): archive cote 134-6, site page 110 (PDF page 111), typescript page 428.
% typescript. Draft compared: ps5.tex:7223-7308.
% AI-transcribed by Claude Opus 5.5 (claude-opus-5-5, adaptive thinking, effort high), Message Batches API, batch ps3 (msgbatch_01XdFmuVVBAcgc7EFxdtCqsj), 2026-11-12; not yet checked by eye.
\scanpage{134-6/110}
\mspage{428}

\uline{Hom}'s componentwise  (and we have the similar formula of course for the
Hom's and RHom\unsure{'}s  non underlined). \xout{Th}

Taking  L. = $\underline{O}_{\hand{\mathcal{X}}}$ (or any resolution of $\underline{O}_{\hand{\mathcal{X}}}$) ,  the RHom invariant
(\hand{*})  reduces to the total derived functor  RI  of the sections functor
\[ \hand{(i)} \qquad \mathrm{RHom}\xout{xx}_{\hand{\underline{O}_{\mathcal{X}}}}(\underline{O}_{\hand{\mathcal{X}}},F^\bullet)\hand{\simeq} \mathrm{R}\over{I}{\Gamma}_{\hand{\mathcal{X}}}(F^\bullet) \]
(whereas  $\mathrm{R}\underline{\mathrm{Hom}}_{\underline{O}_{\hand{\mathcal{X}}}}(\underline{O}_{\hand{\mathcal{X}}},F^\bullet)\hand{\simeq}F^\bullet$  of course),  which in turn allows to give
the following familiar expression of RHom in terms of R\uline{Hom}  :
\[ (\deleted{\illegible{1}})\inserted{(ii)} \qquad \mathrm{RHom}_{\underline{O}_{\hand{\mathcal{X}}}}(L.,F^\bullet)\hand{\simeq} \mathrm{R}\over{I}{\Gamma}_{\hand{\mathcal{X}}}(\mathrm{R}\underline{\mathrm{Hom}}_{\underline{O}_{\hand{\mathcal{X}}}}(L.,F^\bullet)) \qquad , \]
\xout{expressing}\inserted[typed]{coming from} the similar isomorphism  $\mathrm{Hom}_{\underline{O}_{\hand{\mathcal{X}}}}\hand{\simeq}\over{I}{\Gamma}_{\hand{\mathcal{X}}}\underline{\mathrm{Hom}}_{\underline{O}}$ .  All this is
standard cohomology formalism\hand{,} valid on an arbitrary \xout{topos} ringed topos
$(\hand{\mathcal{X}},\underline{O}_{\hand{\mathcal{X}}})$ ,  except for the possibility of computing RHom and R\uline{Hom}  by \xout{takin}
taking \uline{projective resolutions} \uline{of the first argument} (rather than injecti\xout{x}
ve ones of the second),  which is special to the case  when $\hand{\mathcal{X}}$ = A\^{} ,
to which we'll now restrict.

Let now A be a fixed small category, k a fixed commutative ring,
\hand{$\mathcal{M}$} a k-additive category,  stable under small inverse limits.  We want to
define a total derived functor of the functor
\[ (109) \qquad (L,F)\hand{\mapsto}\xout{\illegible{1}}\mathrm{Hom}_k(L,F) \ : \ A^{\wedge}_k \ \mathrm{x} \ A^{\wedge}_{\hand{\mathcal{M}}}\hand{\longrightarrow}\hand{\mathcal{M}}, \]
which should be a functor
\[ (1\over{09}{10}) \qquad (L.,F\unsure{^\bullet})\hand{\mapsto}\mathrm{RHom}_k(L.,F^\bullet) \ : \ D^-(A^{\wedge}_k) \ \mathrm{x} \ D^+(A^{\wedge}_{\hand{\mathcal{M}}})\hand{\longrightarrow} D^+(\hand{\mathcal{M}}\deleted{\illegible{1}}\hand{)}\hand{,} \ \deleted{\illegible{1}} \]
and similarly
\[ (11\over{0}{1}) \qquad (L.,F^\bullet)\hand{\mapsto}\mathrm{R}\underline{\mathrm{Hom}}_k(L.,F^\bullet) \ : \ D^-(A^{\wedge}_k) \ \mathrm{x} \ D^+(A^{\wedge}_{\hand{\mathcal{M}}})\hand{\longrightarrow} D^+(A^{\wedge}_{\hand{\mathcal{M}}}) \quad . \]
For this,  in order for $D^+(A^{\wedge}_{\hand{\mathcal{M}}})$  to be defined,  we better assume \hand{$\mathcal{M}$} to be
an \uline{abelian} category,  hence $A^{\wedge}_{\hand{\mathcal{M}}}$ is abelian too.  \xout{\unsure{With the}} Of course, we'll
write
\[ (11\over{1}{2}) \qquad \mathrm{Ext}^i_k(L\hand{.},F\hand{^\bullet}) = H^i(\mathrm{RHom}_k(L.,F^\bullet)) \ , \ \underline{\mathrm{Ext}}^i_k(L.,F^\bullet) = H^i(\mathrm{R}\underline{\mathrm{Hom}}_k(\hand{\unsure{F_\bullet,L^\bullet}})), \]
these global and local Ext$^{\mathrm{i}}$  may be viewed as  ``external'' Ext$^{\mathrm{i}}$'s ,  as
contrarily to the familiar case,  the \inserted[typed]{components of the} two arguments L. and K$^\bullet$  are not
in the same category  -  just as the Hom$_{\mathrm{k}}$ in (109)  and the corresponding
\[ (109') \qquad \underline{\mathrm{Hom}}_k \ : \ A^{\wedge}_k \ \mathrm{x} \ A^{\wedge}_{\hand{\mathcal{M}}}\hand{\to} A^{\wedge}_{\hand{\mathcal{M}}} \]
has arguments in \inserted{the} two different categories. \hand{$A^{\wedge}_k$ and $A^{\wedge}_{\mathcal{M}}$ .}

\margin{110}
